\section{Feature Engineering and Data Decisions}
The feature engineering follows the order in the handwritten project plan: reduce the data, define the target, select features, encode categorical information, decide about scaling and PCA, then add lag features after the first modelling results showed missing temporal memory.

\begin{table*}[t]
\caption{Main preprocessing and representation decisions.}
\centering
\scriptsize
\begin{tabularx}{\textwidth}{p{0.17\textwidth}Y Y}
\toprule
Decision & Implementation & Justification \\
\midrule
Station-day aggregation & Group trips by \code{start\_station\_id} and date, then merge daily weather. & Gives one consistent observation unit and matches the daily weather granularity. \\
\addlinespace
Feature choice & Keep station identity, selected weather variables, calendar/time variables, and later lag features. & These groups directly match the demand mechanism: location, conditions, seasonality, and recent usage. \\
\addlinespace
Human labels separated & Store station names in a mapping file instead of using names as predictors. & Names help interpretation but duplicate station identity and are not numeric model inputs. \\
\addlinespace
Date transformation & Replace raw date strings with \code{month}, \code{weekday}, \code{year}, and \code{time\_idx}. & Preserves calendar and order information in numeric form. \\
\addlinespace
Leakage control & Remove member/casual and bike-type count decompositions from input features. & These variables are components of the target and would make prediction unrealistically easy. \\
\addlinespace
Correlation reduction & Use Fig.~\ref{fig:corr} to remove near-duplicates such as \code{solarenergy}/\code{solarradiation} and \code{feelslike}/\code{temp}. & Reduces redundant predictors and improves interpretability, especially for linear models. \\
\addlinespace
No PCA & Keep original reduced predictors instead of principal components. & The feature set is already small; PCA would hide the physical meaning of weather, time, and station effects. \\
\addlinespace
One-hot station encoding & Encode \code{start\_station\_id} as binary station columns. & Station IDs are nominal categories, not ordered quantities. One-hot encoding creates one binary feature per category and avoids false numeric distances \cite{sklearn_ohe}. \\
\addlinespace
Selective scaling & Standardize only Linear, Ridge, Lasso, KNN, and Neural Network inputs. & These models use coefficients, penalties, distances, or gradient optimization. Tree models split by thresholds and do not require scaling. \\
\addlinespace
Exact lag features & Add \code{total\_rentals\_lag\_1} and \code{total\_rentals\_lag\_7} per station. & The first no-lag results showed that weather and calendar alone miss temporal dependence. \\
\addlinespace
Fair lag/no-lag comparison & Filter both variants to the same station-day keys. & Prevents a false improvement caused by different rows rather than different features. \\
\bottomrule
\end{tabularx}
\end{table*}

\subsection{Station Identity and One-Hot Encoding}
The first implementation treated \code{start\_station\_id} as a numeric column. This was practical but conceptually wrong: station 31623 is not larger than station 31245 in a meaningful modelling sense. The final version therefore uses one-hot encoding. This is especially important for Linear, Ridge, Lasso, KNN, and Neural Network models because numeric station IDs would create artificial distances or coefficients. Tree models can sometimes split numeric IDs, but one-hot encoding gives a cleaner and comparable representation across all model families.

\subsection{Scaling and PCA}
Scaling is not a universal preprocessing step. It is used only where the model logic needs it. KNN depends on distances, regularized linear models depend on coefficient penalties, and Neural Networks depend on gradient-based optimization. Decision Tree, Random Forest, and Gradient Boosting do not need scaling because their splits are threshold-based. PCA was rejected because the reduced feature set is already manageable and the paper needs interpretable decisions. A principal component may be statistically useful, but it would no longer clearly mean temperature, weekday, station, or lagged demand.

\subsection{Lag Features}
Lag features are computed station-wise, not globally. A global shift would mix demand from different stations and create wrong information. The correct lag for a station-day is the same station's demand one day before or seven days before. Rows without exact lag history are removed. The same rows are also removed from the no-lag variant so that comparisons remain fair. This step directly addresses the main weakness of the preliminary no-lag models: they used current weather and calendar information, but no memory of recent station demand.

\subsection{Why Representation Was Prioritized Over More Algorithms}
Adding more algorithms would not solve a weak data representation. The early results showed that nonlinear models were better than linear models, but the largest conceptual weakness was missing temporal memory. Therefore, the final development focused on station encoding and lag features before interpreting the final ranking. This follows a data-centric approach: improve the input representation first, then compare models under one consistent protocol.
